\section{Moments of the Migdal spectrum}\label{app:D.34}

This section derives the moments of \cref{sec:12.5}: the inverse force moments, the trace rule, the cross and third moments, and the Dirac contact density. The notation is that of \cref{sec:8.6,sec:12.1,sec:12.5}: $|0\rangle$ and $|a\rangle$ are the ground and excited states of the clamped electrons, $\Delta_a = E_a - E_0$, and $\vect{F}_{a0} = \langle a|\nabla_{\vect{R}_\kappa}V_{eN}|0\rangle$.

\subsection{Moments from the Migdal entry}\label{app:D.34-moments}

At leading order the kick of a nucleus reaches the excited surface $a$ with the amplitude \cref{eq:D.28-boost}, $-\ii\,\vect{q}\cdot\vect{d}_{a0}/(M_\kappa\Delta_a)$ times a nuclear overlap. Summing the squared modulus over the complete set of vibrational states of that surface removes the overlap, and the electronic energy of the final state is $\Delta_a$ in the Condon approximation with vertical energies (\cref{sec:12.1-boost}). With $\vect{d}_{a0} = -\vect{F}_{a0}/\Delta_a$ from \cref{eq:5.1-dcoup},
\begin{equation*}
  \frac{\dd P_{\mathrm{Mig}}}{\dd\omega_e} = \sum_{a\neq0}\frac{|\vect{q}\cdot\vect{F}_{a0}|^2}{M_\kappa^2\Delta_a^4}\,\delta(\omega_e - \Delta_a)
  = \frac{q^2}{M_\kappa^2\,\omega_e^4}\,\frac{\dd\sigma_\kappa}{\dd\Delta}\bigg|_{\Delta = \omega_e} ,
\end{equation*}
which is \cref{eq:12.5-spectrum}. The same function is the electronic marginal of the Migdal entry: integrating \cref{eq:12.1-Skk} over $\omega_l$ with the incoherent lattice factor, whose weight is $2\pi\nd_\kappa$ at every momentum (the coherent zero-phonon process for $q \gtrsim 2q_{\mathrm{BZ}}$ on zone average; at $q \ll q_{\mathrm{BZ}}$ that process obeys \cref{eq:8.3-boostward}, \cref{app:D.33-soft}), and summing the electronic factor over final states with \cref{eq:12.1-force}, gives $2\pi\nd_\kappa$ times this expression (\cref{app:D.33-square}). The $k$th moment is $(q^2/M_\kappa^2)\int\Delta^{k-4}\dd\sigma_\kappa = q^2\mathcal{I}_{4-k}/M_\kappa^2$. For $k = 1$, the projection of \cref{eq:D.28-dM} on $\hat{\vect{q}}$ is $2\sum_a|\hat{\vect{q}}\cdot\vect{F}_{a0}|^2/\Delta_a^3 = 2\mathcal{I}_3$, so that $E^{\mathrm{Mig}}_1 = q^2\delta M_\kappa/(2M_\kappa^2) = E_R\,\delta M_\kappa/M_\kappa$. For $k = 2$, the diagonal correction of Born and Huang in the effective Hamiltonian of \cref{app:D.28} is $\sum_a|\vect{d}_{a0}|^2/2M_\kappa$; its part along $\hat{\vect{q}}$ is $E_{\mathrm{BH}} = \mathcal{I}_2/2M_\kappa$, and $E^{\mathrm{Mig}}_2 = q^2\mathcal{I}_2/M_\kappa^2 = 4E_RE_{\mathrm{BH}}$.

For the valence electrons in the soft limit, \cref{eq:12.3-soft} and \cref{eq:12.5-spectrum} describe the same probability, so
\begin{equation*}
  \frac{\dd\sigma_\kappa}{\dd\omega} = \frac{M_\kappa^2\omega^4}{q^2}\,\frac{\dd P_{\mathrm{Mig}}}{\dd\omega}
  = 4\alpha\int\frac{\dd^3p}{(2\pi)^3}\;Z_{\mathrm{ion}}(p)^2\,\frac{(\hat{\vect{q}}\cdot\vect{p})^2}{p^2}\,L(\vect{p},\omega) ,
\end{equation*}
with $\vect{v}_N = \vect{q}/M_\kappa$. Then $\delta M^{\mathrm{val}}_\kappa = 2\int\omega^{-3}\dd\sigma_\kappa$, and for a direction-averaged loss function the angular average of $(\hat{\vect{q}}\cdot\hat{\vect{p}})^2$ is $1/3$, which gives \cref{eq:12.5-dMval}.

\subsection{Dynamical force constant}\label{app:D.34-dynamic}

Displace nucleus $\kappa$ by $u\,\hat{\vect{q}}\cos\omega t$ at a frequency below the gap. The electrons respond linearly, and the force they exert back on the nucleus is $-\Phi^{\mathrm{el}}_{\kappa\kappa}(\omega)u$, with the static term $\langle\partial^2_{\hat q}V_{eN}\rangle$ and the retarded response of the operator $\partial_{\hat q}V_{eN}$ to itself,
\begin{equation*}
  \Phi^{\mathrm{el}}_{\kappa\kappa}(\omega) = \langle\partial^2_{\hat q}V_{eN}\rangle - \sum_{a\neq0}|\hat{\vect{q}}\cdot\vect{F}_{a0}|^2\Bigl[\frac{1}{\Delta_a - \omega} + \frac{1}{\Delta_a + \omega}\Bigr]
  = \langle\partial^2_{\hat q}V_{eN}\rangle - 2\int\dd\sigma_\kappa(\Delta)\,\frac{\Delta}{\Delta^2 - \omega^2} .
\end{equation*}
For $\omega < \Delta$ on the whole support, $\Delta/(\Delta^2 - \omega^2) = \Delta^{-1} + \omega^2\Delta^{-3} + O(\omega^4\Delta^{-5})$, and
\begin{equation*}
  \Phi^{\mathrm{el}}_{\kappa\kappa}(\omega) = \bigl[\langle\partial^2_{\hat q}V_{eN}\rangle - 2\mathcal{I}_1\bigr] - \omega^2\,2\mathcal{I}_3 + O(\omega^4) = \Phi^{\mathrm{el}}_{\kappa\kappa}(0) - \omega^2\delta M_\kappa + O(\omega^4) .
\end{equation*}
The equation of motion $M_\kappa\ddot u = -\Phi u$ at frequency $\omega$ becomes $[\Phi^{\mathrm{el}}_{\kappa\kappa}(0) - \omega^2(M_\kappa + \delta M_\kappa)]u = 0$ to this order. At imaginary frequency the sign of the $\omega^2$ term changes, $\Phi^{\mathrm{el}}_{\kappa\kappa}(\ii\omega) = \Phi^{\mathrm{el}}_{\kappa\kappa}(0) + \omega^2\delta M_\kappa + O(\omega^4)$, which is the form \cref{app:E.18} fits.

\subsection{The trace rule}\label{app:D.34-trace}

Write $V_{eN} = \sum_{i,J}v_J(\vect{r}_i - \vect{R}_J)$, with $i$ the electrons and $J$ the nuclei. Since $\nabla_{\vect{R}_J}v_J(\vect{r}_i - \vect{R}_J) = -\nabla_{\vect{r}_i}v_J(\vect{r}_i - \vect{R}_J)$, the force summed over all nuclei is minus the force on all electrons, $\sum_J\nabla_{\vect{R}_J}V_{eN} = -\sum_i\nabla_{\vect{r}_i}V_{eN}$. The electron--electron interaction depends on differences of electron positions, so $\sum_i\nabla_{\vect{r}_i}V_{ee} = 0$, and with the total electronic momentum $\vect{P}_e = -\ii\sum_i\nabla_{\vect{r}_i}$ and the clamped Hamiltonian $H_{\mathrm{el}}$,
\begin{gather*}
  [\vect{P}_e, H_{\mathrm{el}}] = -\ii\sum_i\nabla_{\vect{r}_i}V_{eN} , \\
  \sum_J\langle a|\nabla_{\vect{R}_J}V_{eN}|0\rangle = -\ii\,\langle a|[\vect{P}_e, H_{\mathrm{el}}]|0\rangle = \ii\Delta_a\,\langle a|\vect{P}_e|0\rangle .
\end{gather*}
Take a finite piece of an insulator, so that the total position $\vect{X}_e = \sum_i\vect{r}_i$ is a bounded operator on its states, and use $\vect{P}_e = \ii m_e[H_{\mathrm{el}}, \vect{X}_e]$, whose matrix element is $\ii m_e\Delta_a\langle a|\vect{X}_e|0\rangle$. Summing \cref{eq:D.28-dM} over both nuclear indices,
\begin{equation*}
  \sum_{I,J}\hat{\vect{q}}\cdot\delta M_{IJ}\cdot\hat{\vect{q}} = 2\sum_{a\neq0}\frac{\Delta_a^2|\hat{\vect{q}}\cdot\vect{P}_{a0}|^2}{\Delta_a^3}
  = 2m_e^2\sum_{a\neq0}\Delta_a\,|\langle a|\hat{\vect{q}}\cdot\vect{X}_e|0\rangle|^2 = N_em_e ,
\end{equation*}
by the Thomas--Reiche--Kuhn rule \cref{eq:4.3-TRK} for the $N_e$ electrons of the piece. The result holds for interacting electrons, since only the double commutator of $\vect{X}_e$ with the kinetic energy enters. In the bulk of a crystal with one atom per cell every nucleus is equivalent, so $\sum_J\delta M_{IJ}$ is independent of $I$ and equals $Z_\kappa m_e$. With several atoms per cell the same sum, written with the sublattice displacements of wavevector $\vect{Q} = 0$, is $N_em_e$ per cell. The surface of the piece changes these numbers by a fraction that vanishes as the piece grows. Both forms are \cref{eq:12.5-trace}. The single element $\delta M_{II}$ contains only $\vect{F}^I_{a0}$, not the sum, and no commutator turns it into a ground-state expectation value; this is why \cref{eq:12.5-onsite} leaves it free.

In the one-dimensional model of \cref{sec:12.5-trace}, with independent electrons, the sum over excited states runs over pairs $v \to c$ from the filled band to the empty ones, and $\delta M_{IJ} = 4\sum_{vc}\Real[\langle v|\partial_{X_I}V|c\rangle\langle c|\partial_{X_J}V|v\rangle]/(\varepsilon_c - \varepsilon_v)^3$, the factor 4 being 2 from \cref{eq:D.28-dM} and 2 from spin. The model is solved in plane waves on a supercell of up to 32 wells, which samples the zone of the primitive cell at 32 points. The trace equals $2\,m_e$ per well to $10^{-4}$, and a single well in a large box gives $\delta M = 2\,m_e$, the Thomas--Reiche--Kuhn value.

\subsection{Cross first moment of an atom}\label{app:D.34-cross}

With $\vect{\rho}_e = \vect{r}_e - \vect{X}$, the center of mass $\vect{R}_{\mathrm{cm}}$ and $\vect{q}_e = m_e\vect{q}/m_A$, \cref{eq:12.2-cm} gives $n_\kappa(\vect{q}) = \ee^{\ii\vect{q}\cdot\vect{X}} = \ee^{\ii\vect{q}\cdot\vect{R}_{\mathrm{cm}}}A_\kappa$ with $A_\kappa = \ee^{-\ii\vect{q}_e\cdot\sum_{e'}\vect{\rho}_{e'}}$. Since $\vect{r}_e = \vect{R}_{\mathrm{cm}} + \vect{\rho}_e - (m_e/m_A)\sum_{e'}\vect{\rho}_{e'}$, the electron density is $n_e(\vect{q}) = \ee^{\ii\vect{q}\cdot\vect{R}_{\mathrm{cm}}}A_e$ with $A_e = \sum_e\ee^{\ii\vect{q}\cdot\vect{\rho}_e}A_\kappa$. Both operators move the center of mass by $\vect{q}$, and a final state is a plane wave of the center of mass times an internal state $|f\rangle$ of energy $E_f$, with the total energy transfer $\omega = q^2/2m_A + E_f - E_0$. For a gas of $\nd_\kappa$ atoms per volume at rest,
\begin{equation*}
  S_{e\kappa}(\vect{q},\omega) = 2\pi\nd_\kappa\sum_f\langle0|A_e^\dagger|f\rangle\langle f|A_\kappa|0\rangle\,\delta\bigl(\omega - q^2/2m_A - E_f + E_0\bigr) .
\end{equation*}
The total first moment $\mathscr{M}^1_{e\kappa} = \int(\dd\omega/2\pi)\,\omega[S_{e\kappa}(\vect{q},\omega) + S_{\kappa e}(-\vect{q},\omega)]$ vanishes, \cref{eq:4.3-species}, because $n_e$ commutes with the nuclear kinetic energy and $n_\kappa$ with everything else. The elastic term $f = 0$ sits at $\omega = q^2/2m_A$. For a real ground state the elements at $-\vect{q}$ are the complex conjugates of those at $\vect{q}$, so the two terms of the moment are equal, and the elastic part is
\begin{equation*}
  \mathscr{M}^1_{e\kappa}\big|_{f = 0} = 2\nd_\kappa\,\frac{q^2}{2m_A}\,\Real\bigl[\langle0|A_e|0\rangle^*\langle0|A_\kappa|0\rangle\bigr] = \nd_\kappa\,\frac{q^2}{m_A}\,F_{\mathrm{cr}}(q) .
\end{equation*}
The excited internal states, bound and ionized, carry the rest of the vanishing total, which is \cref{eq:12.5-cross}. No expansion in $q_e$ has been made. To lowest order, $\langle0|A_\kappa|0\rangle = 1 - O(q_e^2\langle\rho^2\rangle)$ and $\langle0|A_e|0\rangle = f(q)$, the form factor of the electron cloud, so that $F_{\mathrm{cr}} = f(q)$ up to those corrections and terms of relative order $m_e/m_A$.

In a crystal at leading order in $\kappa_{\mathrm{BO}}$, \cref{eq:8.6-fsum-ek} gives the electron-following part as $\nd_\kappa q^2F_{\mathrm{cr}}^*/M_\kappa$ with $F_{\mathrm{cr}}(\vect{q}) = (N_\kappa q^2)^{-1}\langle\sum_{I\in\kappa}\ee^{-\ii\vect{q}\cdot\vect{X}_I}\,(-\ii\vect{q})\cdot\nabla_I\bar\rho_e(\vect{q};u)\rangle$, and the electronic sectors carry its negative. At small $q$ the derivative of the adiabatic density along a displacement of nucleus $I$ is $\ii\vect{q}$ times the charge that follows it, $(Z_\kappa\mathbf{1} - Z^*_\kappa)$ at zero macroscopic field (\cref{sec:8.6}), so $F_{\mathrm{cr}} \to \hat{\vect{q}}\cdot(Z_\kappa\mathbf{1} - Z^*_\kappa)\cdot\hat{\vect{q}}$.

\subsection{The third moment}\label{app:D.34-third}

Let the nucleus $\kappa$ be displaced by $u\hat{\vect{q}}$ in the clamped Hamiltonian. Second-order perturbation theory in $u$ gives the curvature of the ground-state energy,
\begin{equation*}
  \partial^2_{\hat q}E_0 = \langle0|\partial^2_{\hat q}V|0\rangle - 2\sum_{a\neq0}\frac{|\hat{\vect{q}}\cdot\vect{F}_{a0}|^2}{\Delta_a} = \langle\partial^2_{\hat q}V\rangle - 2\mathcal{I}_1 ,
\end{equation*}
with $V = V_{eN} + V_{NN}$; the nucleus--nucleus term has no electronic matrix elements and enters only the first term. With \cref{eq:12.5-moments}, $E^{\mathrm{Mig}}_3 = q^2\mathcal{I}_1/M_\kappa^2 = (q^2/2M_\kappa^2)[\langle\partial^2_{\hat q}V\rangle - \partial^2_{\hat q}E_0]$. For a point nucleus, $v_\kappa(\vect{r}) = -Z_\kappa e^2/(4\pi r)$ and $\nabla^2(1/4\pi r) = -\delta^3(\vect{r})$, so $\nabla^2_{\vect{R}}V_{eN} = Z_\kappa e^2\sum_i\delta^3(\vect{r}_i - \vect{R}_\kappa)$, whose expectation value is $Z_\kappa e^2\nd_e^{(\kappa)}$. The nucleus--nucleus potential is harmonic at $\vect{R}_\kappa$, $\nabla^2_{\vect{R}}V_{NN} = 0$. On a site of cubic symmetry the tensor $\langle\partial_\alpha\partial_\beta V\rangle$ is proportional to the unit tensor, so each direction holds a third of the trace, $\langle\partial^2_{\hat q}V\rangle = Z_\kappa e^2\nd_e^{(\kappa)}/3$. Averaging over the zero-point motion of the nuclei and writing $c_3 = (q^2/2M_\kappa^2)\langle\partial^2_{\hat q}E_0\rangle$ gives \cref{eq:12.5-third}. The leading corrections are those of the leading-order boost: relative order $\kappa_{\mathrm{BO}}^2$, and the ratio of the mean phonon energy to the electronic energies that dominate the moment, which are of the order of the $1s$ binding.

For a nucleus of finite size with charge density $\rho_{\mathrm{nuc}}$, normalized to one, the potential is the Coulomb potential folded with $\rho_{\mathrm{nuc}}$, and $\nabla^2_{\vect{R}}V_{eN} = Z_\kappa e^2\sum_i\rho_{\mathrm{nuc}}(\vect{r}_i - \vect{R}_\kappa)$: the contact density becomes $\int\dd^3r\,\nd_e(\vect{r})\rho_{\mathrm{nuc}}(\vect{r} - \vect{R}_\kappa)$. For Dirac electrons the Hamiltonian depends on $\vect{R}_\kappa$ only through $V_{eN}$, the kinetic term $\boldsymbol{\alpha}_{\mathrm{D}}\cdot\vect{p} + \beta_{\mathrm{D}}m_e$ of \cref{app:C.15} being independent of it, so the steps above hold unchanged with the relativistic density.

\subsection{Dirac density at a finite nucleus}\label{app:D.34-dirac}

The $1s$ solution of the Dirac equation \cref{eq:C.15-HD} in the Coulomb field of a point charge $Z$ has large and small radial components $g(r) = N r^{\gamma_{\mathrm{D}} - 1}\ee^{-\lambda r}$ and $f(r) = -[(1 - \gamma_{\mathrm{D}})/(1 + \gamma_{\mathrm{D}})]^{1/2}g(r)$, with $\gamma_{\mathrm{D}} = [1 - (Z\alpha)^2]^{1/2}$, energy $E = m_e\gamma_{\mathrm{D}}$ and $\lambda = (m_e^2 - E^2)^{1/2} = Z\alpha m_e = Z/a_0$~\cite{bjorken1964relativistic}. These satisfy the radial equations: inserting the power law, the coefficients of $r^{\gamma_{\mathrm{D}} - 2}$ require $\gamma_{\mathrm{D}}^2 = 1 - (Z\alpha)^2$ and those of $r^{\gamma_{\mathrm{D}} - 1}$ fix $\lambda$ and the ratio $f/g$. The density, summed over both components and divided by $4\pi$ for the angles, is
\begin{equation*}
  \nd(r) = \frac{g^2 + f^2}{4\pi} = \frac{N^2}{4\pi}\,\frac{2}{1 + \gamma_{\mathrm{D}}}\,r^{2\gamma_{\mathrm{D}} - 2}\ee^{-2\lambda r} ,
\end{equation*}
and the normalization $\int_0^\infty\dd r\,r^2(g^2 + f^2) = 1$, with $\int_0^\infty\dd r\,r^{2\gamma_{\mathrm{D}}}\ee^{-2\lambda r} = \Gamma(2\gamma_{\mathrm{D}} + 1)/(2\lambda)^{2\gamma_{\mathrm{D}} + 1}$, gives $\nd(r) = (2\lambda)^3(2\lambda r)^{2\gamma_{\mathrm{D}} - 2}\ee^{-2\lambda r}/[4\pi\,\Gamma(2\gamma_{\mathrm{D}} + 1)]$. The nonrelativistic contact density is $\lambda^3/\pi$, \cref{app:C.15}, so the ratio is $2(2\lambda r)^{2\gamma_{\mathrm{D}} - 2}\ee^{-2\lambda r}/\Gamma(2\gamma_{\mathrm{D}} + 1)$, which equals one at $\gamma_{\mathrm{D}} = 1$ and diverges weakly at $r = 0$ otherwise. Inside a nucleus $2\lambda r \le 2ZR_{\mathrm{nuc}}/a_0$, which is $1.8\times10^{-3}$ for silicon and $5.8\times10^{-3}$ for germanium, and we set $\ee^{-2\lambda r} = 1$ there; this overstates the average by \SI{0.14}{\percent} in silicon and \SI{0.43}{\percent} in germanium. We keep the point-charge orbital, whose shape outside the nucleus the finite size changes only at higher order in $2ZR_{\mathrm{nuc}}/a_0$. Averaging over a uniform sphere of radius $R_{\mathrm{nuc}}$, $(3/R_{\mathrm{nuc}}^3)\int_0^{R_{\mathrm{nuc}}}\dd r\,r^2\,r^{2\gamma_{\mathrm{D}} - 2} = 3R_{\mathrm{nuc}}^{2\gamma_{\mathrm{D}} - 2}/(2\gamma_{\mathrm{D}} + 1)$, gives the first form of \cref{eq:12.5-Rrel}. With $2\gamma_{\mathrm{D}} - 2 = -(Z\alpha)^2 + O((Z\alpha)^4)$, the power is $\exp[(Z\alpha)^2\ln(a_0/2ZR_{\mathrm{nuc}})]$, and the prefactor is $1 + O((Z\alpha)^2)$ without a logarithm; this is the second form, with that non-logarithmic remainder written $c_{\mathrm{rel}}(Z\alpha)^2$. The point-charge orbital inside the nucleus fixes the logarithm but not $c_{\mathrm{rel}}$. The density at the nucleus of every $s$ shell has the same radial dependence $r^{2\gamma_{\mathrm{D}} - 2}$ near the origin, set by the nuclear charge alone, and \cref{sec:12.5-third} applies the $1s$ ratio to all of them.
